\chapter{The async syscall ABI}

This chapter describes the userspace--kernel boundary for asynchronous system
calls. It is the concrete interface that sitas (the shard-per-core userspace
runtime) binds to, and it is the design surface where both projects converge.
Every concept described here is implemented and verified on QEMU AArch64 and
x86-64.

\section{The five operations}

\subsection{submit --- start an async operation}

\texttt{submit(asid, op\_code, buffer) -> Result<CompletionCap, SubmitError>}.
Returns immediately with a \texttt{CompletionCap} naming the in-flight
operation. Ownership of the \texttt{buffer} (an owned \texttt{Vec<u8>})
transfers to the kernel for the operation's duration. Returns
\texttt{SubmitError::WouldBlock} when the capability table is full
(submission backpressure).

\subsection{wait --- the reactor's only sleep}

\texttt{wait(asid, cap) -> Result<(), CapError>}. Blocks the calling thread
until the specified capability reaches terminal completion. Internally
registers the thread's \texttt{Waker} as an observer of the
\texttt{Completion} (which is \texttt{Observable}) and yields.

\subsection{poll --- non-blocking check}

\texttt{poll(asid, cap) -> Result<Option<Completed>, CapError>}. Returns
\texttt{Ok(None)} while in flight; returns \texttt{Ok(Some(Completed \{ 
result, buffer \}))} when terminal. The buffer ownership returns to the
caller. Idempotent after completion (returns \texttt{None} once drained).

\subsection{cancel --- drop-as-cancellation}

\texttt{cancel(asid, cap) -> CancelState}. Requests cancellation of an
in-flight operation. If the operation is still in flight, returns
\texttt{CancelRequested}: the kernel retains any transferred buffer until
it posts a terminal \texttt{Cancelled} completion (deferred reclaim,
mirroring sitas's \texttt{defer\_buffer\_drop}). If the operation already
completed, returns \texttt{AlreadyComplete}.

\subsection{close --- release a slot}

\texttt{close(asid, cap) -> Result<(), CapError>}. Frees a capability-table
slot after its terminal completion was posted or its result was consumed by
\texttt{poll}. Fails with \texttt{NotComplete} if the operation is still in
flight.

\section{The exit-observer completion path}

The ABI's intended completion mechanism is declarative: \textbf{the worker
thread exiting IS the completion event}. \texttt{submit\_worker(asid,
worker\_entry, result)} spawns a worker thread and registers the returned
capability as the worker's exit-observer. The worker performs its work and
returns; the thread's \texttt{Thread::Drop} fires the exit-observer, which
calls \texttt{complete()} and wakes any waiter. The worker never references
the capability explicitly.

This is the exact design described in the kernel's own code
(\texttt{scheduler/threads/mod.rs:63--69}) and is implemented in
\texttt{crates/catten/src/completion/mod.rs} (\texttt{CompletionExitObserver}).

\section{The CQ ring (completion queue)}

A shared-memory io\_uring-style ring for zero-syscall completion delivery. A
single 4 KiB page contains a header (head, tail, capacity, overflow) and a
circular array of entries (16 bytes each: \texttt{cap: u64, result: i64}).

\begin{itemize}
    \item \textbf{Producer (kernel):} \texttt{write(cap, result)} advances the
        head. \texttt{complete()} writes to the ring alongside the observer
        signal.
    \item \textbf{Consumer (userspace):} \texttt{read()} advances the tail,
        returns \texttt{Option<CqEntry>}. Zero syscalls.
    \item \textbf{Overflow:} when the ring is full (head + 1 == tail modulo
        capacity), the write is dropped and the overflow counter is
        incremented (non-lossy detection).
\end{itemize}

The ring is allocated from a physical frame and mapped into the user address
space via \texttt{PageType::UserData} (AP\_EL0 accessible). The same physical
memory is accessible from the kernel via HHDM. Currently demonstrated in
kernel-side self-tests; userspace mapping is proven to work (the real-EL0
test maps a CQ ring alongside code and result pages).

\section{Cancellation and the buffer-reclaim contract}

When userspace drops a \texttt{CompletionCap} without awaiting it (or calls
\texttt{cancel}):
\begin{enumerate}
    \item The kernel marks the operation cancelling but \textbf{retains the
        buffer} --- it may still be the target of DMA or kernel access.
    \item A terminal \texttt{Cancelled} completion is eventually posted to
        the CQ ring, handing the buffer back.
    \item On teardown (address-space exit), outstanding buffers are leaked
        rather than freed while the kernel may touch them (the kernel-side
        mirror of sitas's \texttt{mem::forget}-on-timeout rule).
\end{enumerate}

This is not Unix-specific; it is an ownership contract. The ABI is specified
against sitas's io\_uring buffer-ownership discipline, and sitas's existing
tests (\texttt{charlotte\_abi.rs}) validate it.

\section{Backpressure}

\begin{itemize}
    \item \textbf{Submission backpressure:} the per-AS capability table is
        bounded. \texttt{submit} returns \texttt{SubmitError::WouldBlock} when
        full. The caller must drain completions (freeing slots) before
        retrying.
    \item \textbf{Completion backpressure:} the CQ ring is bounded. When full,
        the kernel does \textbf{not} drop completions; it stops producing new
        ones for that shard and marks the queue overflow-pending.
    \item \textbf{Cross-LP backpressure:} \texttt{IPI\_CMD\_QUEUES} are bounded
        per target LP. \texttt{try\_push\_to} returns \texttt{Err(rpc)} on full
        queue. This matches sitas's bounded \texttt{ShardSender<M>} and the
        cache-coherence analysis in its architecture doc.
\end{itemize}

\endinput
